\documentclass[10pt,a4paper]{amsart}
\usepackage[margin=19mm]{geometry}
\usepackage{amsmath,amssymb,mathtools}
\usepackage[scr=rsfs,cal=euler]{mathalfa}
\usepackage{xfrac}
\usepackage[unicode,hidelinks,bookmarks=false]{hyperref}
\newcommand{\ie}{i.e.\ }
\newcommand{\cf}{cf.\ }
\newcommand{\resp}{respectively\ }
\newcommand{\Subsection}{\subsection}
\providecommand{\nobreakdash}{\nobreak\hbox{-}}
\setlength{\emergencystretch}{2em}
\multlinegap0pt
\usepackage{bm}
\usepackage{bbm}
\usepackage{dsfont}
\usepackage{xspace}
\usepackage{cancel}
\usepackage{mathtools}
\usepackage{amsthm}
\makeatletter
\@ifundefined{theorem}{\newtheorem{theorem}{Theorem}}{}
\@ifundefined{lemma}{\newtheorem{lemma}[theorem]{Lemma}}{}
\makeatother
\newcommand{\R}{\mathbb{R}}
\renewcommand{\bar}[1]{\overline{#1}}
\renewcommand{\tilde}{\widetilde}
\title[High-frequency spectral asymptotics and homogenization for q]{High-frequency spectral asymptotics and homogenization for quasiperiodic operators}
\author{Elena Cherkaev, Bryn Davies, S\'ebastien Guenneau, Clemens Thalhammer, Niklas Wellander}
\date{Source: arXiv:2607.03064v1 (CC BY 4.0)}
\begin{document}
\maketitle

\noindent\emph{Source and licence.} High-frequency spectral asymptotics and homogenization for quasiperiodic operators. Version arXiv:2607.03064v1 (CC BY 4.0), distributed under the Creative Commons Attribution 4.0 International licence, as recorded in the arXiv licence metadata for this exact version and shown on the abstract page https://arxiv.org/abs/2607.03064v1. Full rights notice: licenses/arxiv-2607.03064-CC-BY-4.0.txt.\par\smallskip

\noindent\textit{Editorial scope.} Counted unit: the Lemma whose source label is \texttt{lem: FrozenCoefficients} in the article (source0.tex lines 1316--1331, source label \texttt{lem: FrozenCoefficients}), delivered together with the article's own complete original proof of it (source0.tex lines 1332--1390, one \texttt{proof} environment of 59 lines). No mathematics is written, repaired or re-derived: statement and proof are the article's own text line for line. Required context retained uncounted: lem: PseudoeigenvectorsGeneral (lines 1273--1288). Every definition, display and label of those lines is retained unchanged. Omitted from this package and not redistributed: 1--1272, 1289--1315, 1391--1914. Nothing inside a retained excerpt is omitted or altered: every retained body is delivered exactly as the article's lines are written, so no omission receipt of any kind accompanies this package. Citations: the retained excerpts carry no citation command at all, so no bibliography entry of the article is transcribed and no reference is redistributed. Reference adaptation: none was needed and none was made; every reference inside the delivered unit resolves inside it, so no unresolved reference or citation is produced. Printed slips kept literal: no printed word, symbol, hypothesis or number of the counted statement or of its proof was altered. Layout and class: the article's own class is \texttt{amsart} with packages including csquotes, babel, amssymb, microtype, amsmath, amsfonts, amsthm, graphicx, url, appendix, listings, color, comment, caption; the wrapper is the lane's pinned \texttt{amsart} editorial wrapper at 10pt on A4, which restates the theorem-like environments the retained text needs and adds the article's own macro definitions that the retained text needs (\texttt{\textbackslash R}, \texttt{\textbackslash bar}, \texttt{\textbackslash tilde}), with extra wrapper packages (bm, bbm, dsfont, xspace, cancel, mathtools). The delivered numbering is the unit's own. Assets: no retained span contains a figure environment, a graphics-inclusion command, a PDF-page inclusion, a table or a TikZ picture; no image file is copied into this package, the package declares no asset dependency and no image file is redistributed with it. Licence: version arXiv:2607.03064v1 of this article is distributed under the Creative Commons Attribution 4.0 International licence, as recorded in the arXiv licence metadata for this exact version and shown on the abstract page https://arxiv.org/abs/2607.03064v1. A full rights notice is delivered at licenses/arxiv-2607.03064-CC-BY-4.0.txt. Characters: every retained excerpt is pure ASCII, so the delivered engine, XeLaTeX with its default Latin Modern text font, reports no missing character.
\begin{lemma}\label{lem: PseudoeigenvectorsGeneral}
    Let $(\lambda_\eta, u_\eta)$ be a sequence of eigenvalues and unit eigenvectors for $T_\eta$. Assume that $\lambda_\eta\rightarrow\lambda$ and that $\lambda\notin \sigma_{\text{Boundary}}$. Then there exists a sequence of pseudo-eigenvectors $v_\eta$ such that the following hold
    \begin{itemize}
        \item $\operatorname{supp}(v_\eta)\subset\eta^{-1}\bar{\Omega}$,
        \item $v_\eta\in H^{1}$ is a weak solution in $\mathbb{R}^n$ of 
         \begin{equation} 
            \left\{
                \begin{aligned}
                    - \mathrm{div} \; \left( A\left( \eta \vec{y}, \textbf{P} \vec{y} \right) \nabla v_\eta({\vec y}) \right )& = \lambda v_\eta(\vec{y})+r_\eta \;, \qquad \vec{y} \in \eta^{-1}\bar{\Omega} \\
                    v_\eta(\vec{y})& = 0,\qquad \vec{y}\in\partial\eta^{-1}\Omega,
                \end{aligned}
            \right.
        \end{equation}
        where $r_\eta\in H^{-1}(\R^n)$ is an error term with norm converging to $0$.
    \end{itemize}
\end{lemma}

\begin{lemma}\label{lem: FrozenCoefficients}
    Let $(v_\eta)$ be a sequence of pseudo-eigenvectors obtained from Lemma \ref{lem: PseudoeigenvectorsGeneral}. Then there exists a point $\vec x_0\in \Omega$ and a sequence of pseudo-eigenvectors $\tilde{v}_\eta$ such that
    \begin{itemize}
        \item $\operatorname{supp}(\tilde{v}_\eta)\subset\eta^{-1}\bar{\Omega}$,
        \item $\tilde{v}_\eta\in H^1$ is a weak solution in $\mathbb{R}^n$ of 
         \begin{equation} 
            \left\{
                \begin{aligned}
                    - \mathrm{div} \; \left( A\left( \vec{x}_0, \textbf{P} \vec{y} \right) \nabla \tilde{v}_\eta({\vec y}) \right )& = \lambda \tilde{v}_\eta(\vec{y})+\tilde{r}_\eta \;, \qquad \vec{y} \in \eta^{-1}\bar{\Omega} \\
                    \tilde{v}_\eta(\vec{y})& = 0,\qquad \vec{y}\in\partial\eta^{-1}\Omega,
                \end{aligned}
            \right.
        \end{equation}
        where $\tilde{r}_\eta\in H^{-1}(\mathbb{R}^n)$ is  an error term with norm converging to $0$ as $\eta\to0$.
    \end{itemize}
\end{lemma}

\begin{proof}
    We cover the domain $\bar{\Omega}\eta^{-1}$ by a set of non-overlapping cubes $(Q^\eta_i)_{1\leq i\leq N_\eta}$ of size $[-\frac{\beta_\eta}{2\eta}\,\frac{\beta_\eta}{2\eta}]^n$, where $N_\eta = \mathcal{O}(\beta_\eta^{-n})$. For every $\eta$, denote by $i(\eta)$ the index of the cube where the $L^2$ norm of $v_\eta$ restricted to $Q^\eta_{i(\eta)}$ is maximal. Further, let $\vec{y}^\eta_i$ be the centre of the cube $Q^\eta_i$ and let $\vec{x}^\eta_i = \eta\vec{y}^\eta_i$. We note first that
    \begin{equation}\label{eq: FrozenCoefficients1}
        \lVert v_\eta\mid_{Q^\eta_{i(\eta)}}\rVert > c\beta_\eta^{n/2},
    \end{equation}
    for some $c>0$ small enough. Moreover, since $\vec{x}^\eta_{i(\eta)}\in \Omega$, as $\eta\to0$ there exists a convergent subsequence with some limit $\vec{x}_0\in \bar{\Omega}$. Let $\varphi$ be a smooth bump function satisfying
    \begin{equation}
        \begin{cases}
            0 \leq \varphi \leq 1,\\
            \operatorname{supp}(\varphi)\subset[-1,1]^n,\\
            \varphi=1 \text{ in } [-1/2,1/2]^n.
        \end{cases}
    \end{equation}
    Denote $\varphi_\eta(\cdot): = \varphi\left(\frac{\cdot-\vec{y}^\eta_{i(\eta)}}{\eta^{-1}\beta_\eta}\right)$ and further define
    \begin{equation}
        \tilde{v}_\eta := \frac{\varphi_\eta v_\eta}{\lVert \varphi_\eta v_\eta\rVert}.
    \end{equation}
    It remains to check whether $\tilde{v}_\eta$ makes a good sequence of pseudo-eigenvectors for the fixed-variable operator. First, we bound the error occurred by freezing the coefficients: 
    \begin{equation}\label{eq: FrozenCoefficients2}
        \begin{aligned}
            \langle T_{\vec{x}_0}\tilde{v}_\eta,v \rangle_{H^{-1},H^{1}}&=\langle T_{\eta}\tilde{v}_\eta,v \rangle_{H^{-1},H^{1}}\\
            &+ \underbrace{\lVert \varphi_\eta v_\eta\rVert^{-1}\int_{\eta^{-1}\Omega}\left(A\left( \vec{x}_0, \textbf{P}\vec{y}\right)-A\left( \eta\vec{y}, \textbf{P}\vec{y}\right)\right)\nabla  (\varphi_\eta(\vec{y})v_\eta(\vec{y})) \cdot \nabla \overline{v}(\vec{y})dy}_{A_\eta}.    
        \end{aligned}
    \end{equation}
    Thus, we require an upper bound on $A_\eta$
    \begin{equation}\label{eq: FrozenCoefficients3}
        \begin{aligned}
            A_\eta&:=\lVert \varphi_\eta v_\eta\rVert^{-1}\int_{\eta^{-1}\Omega}\left(A\left( \vec{x}_0, \textbf{P}\vec{y}\right)-A\left( \eta\vec{y}, \textbf{P}\vec{y}\right)\right)\nabla  (\varphi_\eta(\vec{y})v_\eta(\vec{y})) \cdot \nabla \overline{v}(\vec{y})dy\\
            &\leq \lVert \varphi_\eta v_\eta\rVert^{-1}[A]_{\mathcal{C}^{0,\alpha}}(\beta_\eta + \lvert \vec{x}_{i(\eta)}^\eta-\vec{x}_0\rvert)^\alpha(\lVert \nabla \varphi_\eta \rVert_{L^\infty}\lVert v_\eta\rVert + \lVert \varphi_\eta \nabla v_\eta\rVert)\lVert\nabla  v\rVert.
        \end{aligned}
    \end{equation}
    For the first term, we note simply that $\varphi$ can be chosen such that $\lvert \nabla \varphi_\eta\rvert \leq C\eta\beta_\eta^{-1}$. To bound the second term, we make use of Caccioppoli's inequality to obtain 
    \begin{equation}
        \lVert \varphi_\eta \nabla v_\eta \rVert^2 \leq C(1 + \lvert \lambda\rvert) \lVert \varphi_\eta v_\eta\rVert^2 + C\lVert r_\eta\rVert_{H^{-1}}^2.   
    \end{equation}
    Applying these bounds to Equation \eqref{eq: FrozenCoefficients3} and making use of \eqref{eq: FrozenCoefficients1}, we obtain 

    \begin{equation}
        A_\eta \leq C[A]_{\mathcal{C}^{0,\alpha}}(\beta_\eta + \lvert \vec{x}_{i(\eta)}^\eta-\vec{x}_0\rvert)^\alpha(\eta\beta_\eta^{-(n+2)/2} + (1 + \lvert \lambda\rvert)^{1/2}+ \beta_\eta^{-n/2}\lVert r_\eta\rVert_{H^{-1}})\lVert\nabla  v\rVert.
    \end{equation}
    In a next step, we examine $\langle T_{\eta}\tilde{v}_\eta,v \rangle_{H^{-1},H^{1}}$
    \begin{equation}
        \begin{aligned}
            &\langle T_{\eta}\tilde{v}_\eta,v \rangle_{H^{-1},H^{1}}=\lVert \varphi_\eta v_\eta\rVert^{-1}\int_{\eta^{-1}\Omega}A\left( \eta\vec{y}, \textbf{P}\vec{y}\right)\nabla  (\varphi_\eta(\vec{y})v_\eta(\vec{y})) \cdot \nabla \overline{v}(\vec{y})dy\\
            &\quad =\lVert \varphi_\eta v_\eta\rVert^{-1}\langle T_{\eta}v_\eta,
            \varphi_\eta v \rangle_{H^{-1},H^{1}}\\
            & + \underbrace{\lVert \varphi_\eta v_\eta\rVert^{-1}\int_{\eta^{-1}\Omega}A\left( \eta\vec{y}, \textbf{P}\vec{y}\right) \left((\nabla  \varphi_\eta(\vec{y}))v_\eta(\vec{y})) \cdot \nabla \overline{v}(\vec{y}) - \nabla v_\eta \cdot (\nabla \varphi_\eta)(\vec{y})\overline{v}(\vec{y})\right)d\vec{y}}_{B_\eta}.
        \end{aligned}
    \end{equation}
    Since every term in $B_\eta$ includes a factor of $\nabla \varphi_\eta$, and we can again exploit the choice that $\lvert \nabla \varphi_\eta\rvert \leq C\eta\beta_\eta^{-1}$, we have the bound
    \begin{equation}
        B_\eta\leq  C\eta\beta_\eta^{-1}\lVert \varphi_\eta v_\eta\rVert^{-1} \lVert A\rVert_{L^\infty}((1+\lvert \lambda \rvert)^{1/2}\lVert \varphi_\eta v_\eta\rVert + \lVert r_\eta\rVert_{H^{-1}})\lVert v\rVert_{H^1}.
    \end{equation}
    Putting everything together, we have
    \begin{equation}
        \langle T_{\vec{x}_0}\tilde{v}_\eta,v \rangle_{H^{-1},H^{1}} = \lambda \langle \tilde{v}_\eta,v \rangle + \lVert \varphi_\eta v_\eta\rVert^{-1}\langle r_\eta,v\rangle_{H^{-1},H^1}+ A_\eta + B_\eta.
    \end{equation}
    For $\beta_\eta = \max(\eta^{2/(n+3)}, \lVert r_\eta\rVert_{H^{-1}}^{1/n})$ the claim follows.
\end{proof}

\end{document}
